\section{Validação dos canais selecionados}
\label{sec:validacao}

A validação será conduzida como um piloto de quatro semanas em uma região de Belo Horizonte, com pelo menos duas ONGs e duas clínicas parceiras. O objetivo é verificar se cada canal gera descoberta, ativação, uso e continuidade, sem confundir interesse inicial com valor entregue.

\begin{table}[htbp]
\centering\small
\caption{Plano inicial de testes}
\label{tab:testes}
\begin{tabularx}{\textwidth}{p{0.18\textwidth}p{0.27\textwidth}Y}
\toprule
\textbf{Frente} & \textbf{Teste} & \textbf{Sinal de aprendizado}\\
\midrule
Comunicação & Publicar duas mensagens equivalentes em grupos de bairro e em Instagram/Facebook, com chamadas distintas para cadastro. & Comparar alcance, cliques e cadastros qualificados por origem.\\
Comunicação & Entregar cartões/QR codes diferentes em duas ONGs e duas clínicas. & Identificar quais parceiros geram mais ativações e quais perfis chegam por indicação.\\
Distribuição & Conduzir onboarding de um grupo pelo webapp e de outro pelo WhatsApp. & Comparar conclusão do cadastro, tempo até o primeiro lembrete e necessidade de ajuda humana.\\
Distribuição & Oferecer uma agenda piloto com número limitado de horários sociais. & Medir procura, agendamentos, comparecimento e uso de créditos.\\
Relacionamento & Enviar lembrete de vacinação/vermifugação com janela de resposta e oferecer suporte humano. & Verificar leitura, resposta, ação concluída e percepção de utilidade.\\
Relacionamento & Criar uma comunidade moderada com conteúdo semanal e escuta dos tutores. & Medir participação, dúvidas resolvidas e sugestões convertidas em melhorias.\\
\bottomrule
\end{tabularx}
\end{table}

A coleta deve registrar origem do usuário, categoria de canal, etapa alcançada e feedback. Para preservar confiança, o piloto deve solicitar consentimento, minimizar dados pessoais e manter a distinção entre orientação automatizada e atendimento veterinário.
